\chapter{ĐÁNH GIÁ ĐỊNH LƯỢNG HỆ THỐNG AI VÀ KINH TẾ VẬN HÀNH}

\section{Phương pháp luận và bộ dữ liệu đánh giá thực nghiệm}

Để đánh giá khách quan chất lượng của tầng trí tuệ nhân tạo trong GarageTools AI thay vì chỉ dựa trên cảm quan định tính, nhóm đã xây dựng bộ dữ liệu kiểm định chuẩn hóa gồm \textbf{320 mẫu thử nghiệm} (80 mẫu cho mỗi tác vụ trong 4 tác vụ AI01--AI04), bao gồm:
\begin{itemize}
  \item \textbf{240 kịch bản nghiệp vụ tiêu chuẩn}: Các ca bảo dưỡng định kỳ, mòn má phanh, hỏng máy phát điện, rò rỉ môi chất lạnh điều hòa, lỗi cảm biến oxy, trượt côn và rung giật hộp số trên các dòng xe phổ biến tại Việt Nam (Toyota, Hyundai, Mazda, Ford, Honda, Kia).
  \item \textbf{80 kịch bản biên và gây nhiễu (Edge \& Stress Cases)}: Mô tả triệu chứng mơ hồ, yêu cầu phụ tùng đã hết hàng trong kho ($\text{AvailableQty} = 0$), hoặc chứa câu lệnh tiêm nhiễm nhằm thử thách khả năng bám sát dữ liệu thực (Grounding) của mô hình.
\end{itemize}

Ba thước đo định lượng cốt lõi được sử dụng gồm:
\begin{enumerate}
  \item \textbf{Tỷ lệ hợp lệ cấu trúc (Schema Validation Rate -- SVR)}: Tỷ lệ phần trăm phản hồi từ LLM vượt qua bộ kiểm định \texttt{Pydantic v2} ngay ở lần phân tích đầu tiên mà không gặp lỗi cú pháp JSON hay thiếu trường bắt buộc.
  \item \textbf{Điểm bám sát dữ liệu gốc (Grounding Score -- GS)}: Tỷ lệ các mã phụ tùng (\texttt{sku}), đơn giá bán và lịch sử ODO trong câu trả lời của AI khớp chính xác 100\% với dữ liệu tồn kho và hồ sơ xe đã nạp vào vùng ngữ cảnh RAG.
  \item \textbf{Tỷ lệ ảo giác (Hallucination Rate -- HR)}: Tỷ lệ các phản hồi tự bịa đặt mã phụ tùng không có trong bảng \texttt{parts}, sai lệch số học tồn kho hoặc đưa ra kết luận vượt quá dữ liệu đầu vào ($\text{HR} = 100\% - \text{GS}$).
\end{enumerate}

\section{Đánh giá sự tiến hóa qua 3 vòng tối ưu hóa Prompt}

Bảng~\ref{tab:g-prompt-eval} ghi nhận sự cải thiện rõ rệt của các chỉ số chất lượng qua 3 vòng tối ưu hóa Prompt (\texttt{v1.0} $\rightarrow$ \texttt{v2.0} $\rightarrow$ \texttt{v3.2}) trên tập dữ liệu thử nghiệm.

\begin{table}[H]
\centering
\footnotesize
\renewcommand{\arraystretch}{1.02}
\caption{Kết quả đánh giá định lượng qua 3 vòng lặp tối ưu hóa Prompt (N = 320 mẫu)}
\label{tab:g-prompt-eval}
\begin{tabularx}{\textwidth}{|L{2.8cm}|C{2.0cm}|C{2.1cm}|C{2.0cm}|Y|}
\hline
\thead{Phiên bản Prompt} & \thead{Schema Valid} & \thead{Grounding Score} & \thead{Tỷ lệ ảo giác} & \thead{Nhận xét kỹ thuật từ thực nghiệm}\\
\hline
Vòng 1: \texttt{v1.0} (Zero-shot cơ bản) & 35.0\% & 58.2\% & 41.8\% & Thường xuyên trả về văn bản tự do ngoài khung JSON; tự bịa mã phụ tùng ngoại lai và đơn giá ước lượng bằng USD. \\
\hline
Vòng 2: \texttt{v2.0} (Role + JSON Schema) & 88.5\% & 81.0\% & 19.0\% & Định dạng JSON ổn định hơn nhưng đôi lúc vẫn gợi ý phụ tùng đã hết tồn kho khả dụng mà không bật cảnh báo. \\
\hline
Vòng 3: \texttt{v3.2} (RAG SKU + Guardrails) & \textbf{100.0\%} & \textbf{96.4\%} & \textbf{3.6\%} & Khớp 100\% hợp đồng dữ liệu Pydantic v2; chỉ chọn mã SKU có thật trong kho và tự động gắn cờ khi thiếu tồn kho. \\
\hline
\end{tabularx}
\end{table}

\section{Kết quả thực nghiệm định lượng theo từng tác vụ AI}

Với cấu hình Prompt phiên bản \texttt{v3.2} kết hợp cơ chế hậu kiểm đối chiếu CSDL và Smart Fallback, kết quả đo đạc chi tiết trên từng tác vụ AI01--AI04 được tổng hợp tại Bảng~\ref{tab:g-task-metrics}.

\begin{table}[H]
\centering
\footnotesize
\renewcommand{\arraystretch}{1.02}
\caption{Chỉ số chất lượng thực nghiệm phân theo 4 tác vụ AI cốt lõi (Phiên bản Prompt v3.2)}
\label{tab:g-task-metrics}
\begin{tabularx}{\textwidth}{|L{3.5cm}|C{1.6cm}|C{1.9cm}|C{1.9cm}|Y|}
\hline
\thead{Tác vụ AI} & \thead{Số mẫu} & \thead{Schema Valid} & \thead{Grounding} & \thead{Cơ chế kiểm soát chất lượng đặc thù}\\
\hline
AI01: Chẩn đoán triệu chứng & 80 & 100.0\% & 95.5\% & Đối chiếu lịch sử sửa chữa của mã VIN; phân loại mức độ khẩn cấp và yêu cầu KTV kiểm tra thực tế. \\
\hline
AI02: Đề xuất báo giá \& SKU & 80 & 100.0\% & 97.2\% & Hậu kiểm mã \texttt{sku} trực tiếp với bảng \texttt{parts}; tự động gắn cờ \texttt{stock\_warning} nếu $\text{Available} < \text{Qty}$. \\
\hline
AI03: Tóm tắt tiến độ RO & 80 & 100.0\% & 96.8\% & Chỉ tổng hợp từ các hạng mục có thật trong \texttt{ro\_line\_items} và trạng thái nghiệm thu \texttt{qc\_inspections}. \\
\hline
AI04: Cảnh báo bảo dưỡng & 80 & 100.0\% & 96.1\% & Nội suy mốc bảo dưỡng kế tiếp dựa trên số ODO hiện tại và định mức chu kỳ kỹ thuật chuẩn. \\
\hline
\textbf{Trung bình toàn hệ thống} & \textbf{320} & \textbf{100.0\%} & \textbf{96.4\%} & \textbf{0\% lỗi cấu trúc; 100\% qua kiểm duyệt HITL trước khi áp dụng} \\
\hline
\end{tabularx}
\end{table}

\section{Phân tích hiệu năng độ trễ (Latency P50 / P95 / P99)}

Nhóm đã tiến hành đo đạc thời gian phản hồi thực tế (End-to-End Latency) của các nhóm API nghiệp vụ nội bộ và các chế độ gọi AI, thể hiện qua Bảng~\ref{tab:g-latency} và biểu đồ phân bố tại Hình~\ref{fig:g-latency-chart}.

\begin{table}[H]
\centering
\footnotesize
\renewcommand{\arraystretch}{1.02}
\caption{Thống kê phân phối độ trễ phản hồi của hệ thống theo phân vị (P50, P95, P99)}
\label{tab:g-latency}
\begin{tabularx}{\textwidth}{|L{4.6cm}|C{1.8cm}|C{1.8cm}|C{1.8cm}|Y|}
\hline
\thead{Nhóm tác vụ / Chế độ vận hành} & \thead{P50 (Trung vị)} & \thead{P95} & \thead{P99} & \thead{Đánh giá theo ngưỡng NFR03 \& NFR05}\\
\hline
API Nghiệp vụ CRUD \& Tra cứu kho & 8.4 ms & 24.6 ms & 38.2 ms & Đạt vượt mức yêu cầu (Ngưỡng mục tiêu P95 $<$ 45 ms). \\
\hline
Giao dịch Giữ kho nguyên tử (\texttt{BEGIN IMMEDIATE}) & 11.2 ms & 29.5 ms & 43.0 ms & Bảo đảm khóa nguyên tử nhanh, không gây nghẽn cổ chai. \\
\hline
Động cơ AI Smart Fallback (Nội bộ) & 14.5 ms & 22.8 ms & 31.4 ms & Phản hồi tức thời khi đứt kết nối mạng (Mục tiêu $<$ 50 ms). \\
\hline
Groq Cloud LLM (\texttt{gpt-oss-120b}) & 840 ms & 1.650 ms & 2.320 ms & Tốc độ suy luận đám mây rất nhanh, trải nghiệm mượt mà. \\
\hline
Google Gemini (\texttt{gemini-2.5-flash}) & 1.380 ms & 2.740 ms & 3.850 ms & Ổn định, nằm hoàn toàn trong ngưỡng Timeout an toàn (12.000 ms). \\
\hline
\end{tabularx}
\end{table}

\begin{figure}[H]
\centering
\resizebox{0.96\textwidth}{!}{%
\begin{tikzpicture}[font=\sffamily\scriptsize]
  % Legend at top
  \draw[fill=black!18, draw=black, line width=0.6pt] (3.2, 4.35) rectangle (3.6, 4.6);
  \node[anchor=west, font=\sffamily\footnotesize\bfseries] at (3.7, 4.48) {Trung vị P50 (ms)};
  \draw[fill=black!50, draw=black, line width=0.6pt] (7.5, 4.35) rectangle (7.9, 4.6);
  \node[anchor=west, font=\sffamily\footnotesize\bfseries] at (8.0, 4.48) {Phân vị P95 (ms)};

  % Left Panel: Internal APIs & Smart Fallback (0 - 35 ms, scale: 10 ms = 1.0 cm)
  \node[font=\sffamily\footnotesize\bfseries] at (3.3, 3.85) {(a) Nội bộ \& Smart Fallback (0--35 ms)};
  \draw[line width=0.6pt] (0.6, 0) -- (0.6, 3.5);
  \draw[line width=0.6pt] (0.6, 0) -- (6.2, 0);
  \foreach \y/\val in {0/0, 1.0/10, 2.0/20, 3.0/30} {
    \draw[dashed, black!30] (0.6, \y) -- (6.2, \y);
    \node[anchor=east] at (0.5, \y) {\val};
  }
  % Bar group 1: API CRUD (8.4, 24.6)
  \draw[fill=black!18, draw=black, line width=0.6pt] (1.1, 0) rectangle (1.55, 0.84);
  \node[above] at (1.325, 0.84) {8.4};
  \draw[fill=black!50, draw=black, line width=0.6pt] (1.6, 0) rectangle (2.05, 2.46);
  \node[above] at (1.825, 2.46) {24.6};
  \node[below, align=center] at (1.575, -0.05) {API CRUD};

  % Bar group 2: Giữ kho ACID (11.2, 29.5)
  \draw[fill=black!18, draw=black, line width=0.6pt] (2.9, 0) rectangle (3.35, 1.12);
  \node[above] at (3.125, 1.12) {11.2};
  \draw[fill=black!50, draw=black, line width=0.6pt] (3.4, 0) rectangle (3.85, 2.95);
  \node[above] at (3.625, 2.95) {29.5};
  \node[below, align=center] at (3.375, -0.05) {Giữ kho};

  % Bar group 3: AI Fallback (14.5, 22.8)
  \draw[fill=black!18, draw=black, line width=0.6pt] (4.7, 0) rectangle (5.15, 1.45);
  \node[above] at (4.925, 1.45) {14.5};
  \draw[fill=black!50, draw=black, line width=0.6pt] (5.2, 0) rectangle (5.65, 2.28);
  \node[above] at (5.425, 2.28) {22.8};
  \node[below, align=center] at (5.175, -0.05) {Fallback};

  % Right Panel: Cloud LLM (0 - 3000 ms, scale: 1000 ms = 1.0 cm)
  \node[font=\sffamily\footnotesize\bfseries] at (10.7, 3.85) {(b) Cloud LLM qua HTTPS (0--3.000 ms)};
  \draw[line width=0.6pt] (8.0, 0) -- (8.0, 3.5);
  \draw[line width=0.6pt] (8.0, 0) -- (13.2, 0);
  \foreach \y/\val in {0/0, 1.0/1.000, 2.0/2.000, 3.0/3.000} {
    \draw[dashed, black!30] (8.0, \y) -- (13.2, \y);
    \node[anchor=east] at (7.9, \y) {\val};
  }
  % Bar group 4: Groq LLM (840, 1650)
  \draw[fill=black!18, draw=black, line width=0.6pt] (8.9, 0) rectangle (9.45, 0.84);
  \node[above] at (9.175, 0.84) {840};
  \draw[fill=black!50, draw=black, line width=0.6pt] (9.55, 0) rectangle (10.1, 1.65);
  \node[above] at (9.825, 1.65) {1.650};
  \node[below, align=center] at (9.5, -0.05) {Groq LLM};

  % Bar group 5: Gemini Flash (1380, 2740)
  \draw[fill=black!18, draw=black, line width=0.6pt] (11.1, 0) rectangle (11.65, 1.38);
  \node[above] at (11.375, 1.38) {1.380};
  \draw[fill=black!50, draw=black, line width=0.6pt] (11.75, 0) rectangle (12.3, 2.74);
  \node[above] at (12.025, 2.74) {2.740};
  \node[below, align=center] at (11.7, -0.05) {Gemini 2.5 Flash};
\end{tikzpicture}%
}
\caption{Biểu đồ đối sánh độ trễ P50 và P95 (đơn vị: ms) giữa các phân hệ nội bộ và dịch vụ Cloud LLM}
\label{fig:g-latency-chart}
\end{figure}

\section{Phân tích kinh tế vận hành (Token Economics \& Cost Analysis)}

Nhờ kỹ thuật chọn lọc ngữ cảnh RAG (chỉ truyền các phụ tùng liên quan thay vì đổ toàn bộ cơ sở dữ liệu vào Prompt), dung lượng trung bình của mỗi yêu cầu AI được tối ưu hóa ở mức $\approx 620$ Input Tokens và $\approx 280$ Output Tokens. Bảng~\ref{tab:g-cost} phân tích chi phí thực tế trong giai đoạn phát triển và dự phóng chi phí khi vận hành thương mại tại một xưởng dịch vụ quy mô 1.000 lượt gọi AI mỗi tháng.

\begin{table}[H]
\centering
\small
\caption{Phân tích chi phí sử dụng Token và dự phóng kinh tế vận hành cho xưởng dịch vụ}
\label{tab:g-cost}
\begin{tabularx}{\textwidth}{|L{3.8cm}|C{2.2cm}|C{2.4cm}|Y|}
\hline
\thead{Chế độ / Nhà cung cấp AI} & \thead{Token TB / Lượt} & \thead{Chi phí / 1.000 lượt} & \thead{Đánh giá hiệu quả kinh tế vận hành}\\
\hline
Google AI Studio / Groq Free Tier + Smart Fallback & 620 In / 280 Out & \textbf{0 VNĐ} & Chi phí thực tế trong quá trình phát triển và kiểm thử đồ án bằng 0 VNĐ nhờ tận dụng hạn ngạch miễn phí. \\
\hline
Thương mại: Gemini 2.5 Flash / GPT-4o-mini (Pay-as-you-go) & 620 In / 280 Out & $\approx \textbf{2.850 VNĐ}$ ($\sim \$0.11$) & Chi phí vận hành cực thấp: trung bình một xưởng dịch vụ chỉ tốn dưới 3.000 VNĐ cho mỗi 1.000 lần trợ lý AI hỗ trợ lập báo giá. \\
\hline
\end{tabularx}
\end{table}

\section{Kết luận chương}

Chương 7 đã chứng minh định lượng hiệu quả của kiến trúc AI trong GarageTools AI: quá trình tối ưu hóa Prompt 3 vòng kết hợp RAG Grounding đã đưa tỷ lệ hợp lệ Schema đạt 100\%, điểm bám sát dữ liệu thực đạt 96.4\%, độ trễ trung vị dưới 1,4 giây và chi phí vận hành thương mại chưa tới 3.000 VNĐ cho 1.000 lượt yêu cầu.
